% GENERATED FILE. Edit canonical lessons, then run npm run generate:books.
% canonical-source-hash: fnv1a64:6e55037ea939fbaf
% canonical-lessons: MR-C403-khadda, MR-C403-gatirodhak, MR-C403-khadbadit, MR-C403-kalandne, MR-C403-apatne, MR-R403-review-396-403

\chapter{Potholes and Bumps}
\label{ch:mr-potholes-and-bumps}
\hlchaptermodality{403}

\begin{chapteropening}
\textbf{By the end of this chapter:} \emph{I can talk about potholes, speed breakers and a bumpy ride.}

Complete the last lesson of chapter 403: I can talk about potholes, speed breakers and a bumpy ride.
\end{chapteropening}

\section[\texorpdfstring{\mr{खड्डा} (khaḍḍā)}{खड्डा (khaḍḍā)}]{\mr{खड्डा} (khaḍḍā) — a pothole; a pit}

\label{lesson:MR-C403-khadda}



\begin{quote}
Before the new one: say the Marathi for winding, then the Marathi for a rock pinnacle.
\end{quote}

\subsection*{You'll want to know: \mr{खड्डा}}
\textbf{\mr{खड्डा}} — \emph{khaḍḍā} — \textquotedblleft{}a pothole; a pit\textquotedblright{}.

\begin{quote}
\mr{पावसाळ्यात} \mr{रस्त्यावर} \mr{खूप} \mr{खड्डे} \mr{पडतात}.
\end{quote}

\begin{quote}
\emph{pāvsāḷyāt rastyāvar khūp khaḍḍe paḍtāt.}
\end{quote}

\begin{quote}
In the rainy season lots of potholes form on the road.
\end{quote}

Masculine. Plural: \mr{खड्डे}, \emph{khaḍḍe}.

\mr{खड्डा} \mr{पडणे}, a pothole forms; the plural is \mr{खड्डे}.

\begin{cousinweb}
\textbf{\mr{खड्डा}} (\emph{khaḍḍā}) is inherited: CDIAL lists Marathi \emph{khaḍḍā}, \textquotedblleft{}rough hole, pit\textquotedblright{}, under \emph{*khaḍḍa}, \textquotedblleft{}hole, pit\textquotedblright{} (CDIAL 3790).
\end{cousinweb}

\begin{grammarlens}[title={What you've built}]
One more word for getting around.
\end{grammarlens}

\subsection*{Your turn}
\begin{itemize}
\raggedright
  \item \emph{Say it:} \emph{khaḍḍā}
  \item \emph{Say it:} \emph{khaḍḍā}, once more
  \item \emph{Say it:} \emph{pāvsāḷyāt rastyāvar khūp khaḍḍe paḍtāt.} — In the rainy season lots of potholes form on the road.
  \item \emph{Recall:} say the Marathi for winding, then the Marathi for a rock pinnacle, then say \emph{khaḍḍā} again
\end{itemize}

\begin{tcolorbox}[breakable,colback=teal!4,colframe=teal!35!black,title={Before you move on}]
What does \mr{खड्डा} mean? (\textquotedblleft{}A pothole; a pit\textquotedblright{}.) Say it once more.
\end{tcolorbox}
\section[\texorpdfstring{\mr{गतिरोधक} (gatirodhak)}{गतिरोधक (gatirodhak)}]{\mr{गतिरोधक} (gatirodhak) — a speed breaker}

\label{lesson:MR-C403-gatirodhak}



\begin{quote}
Before the new one: say the Marathi for a rock pinnacle, then the Marathi for a pothole.
\end{quote}

\subsection*{You'll want to know: \mr{गतिरोधक}}
\textbf{\mr{गतिरोधक}} — \emph{gatirodhak} — \textquotedblleft{}a speed breaker\textquotedblright{}.

\begin{quote}
\mr{शाळेजवळ} \mr{दोन} \mr{गतिरोधक} \mr{आहेत}.
\end{quote}

\begin{quote}
\emph{śāḷejavaḷ don gatirodhak āhet.}
\end{quote}

\begin{quote}
There are two speed breakers near the school.
\end{quote}

Masculine.

The plural is the same: \mr{दोन} \mr{गतिरोधक}.

\begin{cousinweb}
\textbf{\mr{गतिरोधक}} (\emph{gatirodhak}) is a modern word built from Sanskrit \emph{gati}, \textquotedblleft{}motion\textquotedblright{}, and \emph{rodhaka}, \textquotedblleft{}blocking, checking\textquotedblright{} (Monier-Williams): a thing that holds speed back.
\end{cousinweb}

\begin{grammarlens}[title={What you've built}]
One more word for getting around.
\end{grammarlens}

\subsection*{Your turn}
\begin{itemize}
\raggedright
  \item \emph{Say it:} \emph{gatirodhak}
  \item \emph{Say it:} \emph{gatirodhak}, once more
  \item \emph{Say it:} \emph{śāḷejavaḷ don gatirodhak āhet.} — There are two speed breakers near the school.
  \item \emph{Recall:} say the Marathi for a rock pinnacle, then the Marathi for a pothole, then say \emph{gatirodhak} again
\end{itemize}

\begin{tcolorbox}[breakable,colback=teal!4,colframe=teal!35!black,title={Before you move on}]
What does \mr{गतिरोधक} mean? (\textquotedblleft{}A speed breaker\textquotedblright{}.) Say it once more.
\end{tcolorbox}
\section[\texorpdfstring{\mr{खडबडीत} (khaḍbaḍīt)}{खडबडीत (khaḍbaḍīt)}]{\mr{खडबडीत} (khaḍbaḍīt) — bumpy; rough}

\label{lesson:MR-C403-khadbadit}



\begin{quote}
Before the new one: say the Marathi for a pothole, then the Marathi for a speed breaker.
\end{quote}

\subsection*{You'll want to know: \mr{खडबडीत}}
\textbf{\mr{खडबडीत}} — \emph{khaḍbaḍīt} — \textquotedblleft{}bumpy; rough\textquotedblright{}.

\begin{quote}
\mr{हा} \mr{रस्ता} \mr{फार} \mr{खडबडीत} \mr{आहे}.
\end{quote}

\begin{quote}
\emph{hā rastā phār khaḍbaḍīt āhe.}
\end{quote}

\begin{quote}
This road is very bumpy.
\end{quote}

The opposite is \mr{गुळगुळीत}, smooth; both end in -\mr{ईत} and do not change for gender.

\begin{cousinweb}
\textbf{\mr{खडबडीत}} (\emph{khaḍbaḍīt}) is built on the echo word \emph{khaḍbaḍ}, a rough, knocking noise, with the adjective ending -\mr{ईत}; its deeper origin is not settled.
\end{cousinweb}

\begin{grammarlens}[title={What you've built}]
One more word for getting around.
\end{grammarlens}

\subsection*{Your turn}
\begin{itemize}
\raggedright
  \item \emph{Say it:} \emph{khaḍbaḍīt}
  \item \emph{Say it:} \emph{khaḍbaḍīt}, once more
  \item \emph{Say it:} \emph{hā rastā phār khaḍbaḍīt āhe.} — This road is very bumpy.
  \item \emph{Recall:} say the Marathi for a pothole, then the Marathi for a speed breaker, then say \emph{khaḍbaḍīt} again
\end{itemize}

\begin{tcolorbox}[breakable,colback=teal!4,colframe=teal!35!black,title={Before you move on}]
What does \mr{खडबडीत} mean? (\textquotedblleft{}Bumpy; rough\textquotedblright{}.) Say it once more.
\end{tcolorbox}
\section[\texorpdfstring{\mr{कलंडणे} (kalaṇḍṇe)}{कलंडणे (kalaṇḍṇe)}]{\mr{कलंडणे} (kalaṇḍṇe) — to tip over}

\label{lesson:MR-C403-kalandne}



\begin{quote}
Before the new one: say the Marathi for a speed breaker, then the Marathi for bumpy.
\end{quote}

\subsection*{You'll want to know: \mr{कलंडणे}}
\textbf{\mr{कलंडणे}} — \emph{kalaṇḍṇe} — \textquotedblleft{}to tip over\textquotedblright{}.

\begin{quote}
\mr{वळणावर} \mr{भाजीची} \mr{गाडी} \mr{कलंडली}.
\end{quote}

\begin{quote}
\emph{vaḷaṇāvar bhājīcī gāḍī kalaṇḍlī.}
\end{quote}

\begin{quote}
The vegetable cart tipped over at the bend.
\end{quote}

Past (he …): \mr{तो} \mr{कलंडला}, \emph{to kalaṇḍlā}.

It is something heavy leaning until it falls on its side.

\begin{cousinweb}
\textbf{\mr{कलंडणे}} (\emph{kalaṇḍṇe}) is a Marathi verb for a load or vehicle tipping onto its side; its origin is not settled.
\end{cousinweb}

\begin{grammarlens}[title={What you've built}]
One more word for getting around.
\end{grammarlens}

\subsection*{Your turn}
\begin{itemize}
\raggedright
  \item \emph{Say it:} \emph{kalaṇḍṇe}
  \item \emph{Say it:} \emph{kalaṇḍṇe}, once more
  \item \emph{Say it:} \emph{vaḷaṇāvar bhājīcī gāḍī kalaṇḍlī.} — The vegetable cart tipped over at the bend.
  \item \emph{Recall:} say the Marathi for a speed breaker, then the Marathi for bumpy, then say \emph{kalaṇḍṇe} again
\end{itemize}

\begin{tcolorbox}[breakable,colback=teal!4,colframe=teal!35!black,title={Before you move on}]
What does \mr{कलंडणे} mean? (\textquotedblleft{}To tip over\textquotedblright{}.) Say it once more.
\end{tcolorbox}
\section[\texorpdfstring{\mr{आपटणे} (āpaṭṇe)}{आपटणे (āpaṭṇe)}]{\mr{आपटणे} (āpaṭṇe) — to bang against}

\label{lesson:MR-C403-apatne}



\begin{quote}
Before the new one: say the Marathi for bumpy, then the Marathi for to tip over.
\end{quote}

\subsection*{You'll want to know: \mr{आपटणे}}
\textbf{\mr{आपटणे}} — \emph{āpaṭṇe} — \textquotedblleft{}to bang against\textquotedblright{}.

\begin{quote}
\mr{खड्ड्यामुळे} \mr{माझं} \mr{डोकं} \mr{खिडकीला} \mr{आपटलं}.
\end{quote}

\begin{quote}
\emph{khaḍḍyāmuḷe mājhaṁ ḍokaṁ khiḍkīlā āpaṭlaṁ.}
\end{quote}

\begin{quote}
Because of the pothole my head banged against the window.
\end{quote}

Past (he …): \mr{तो} \mr{आपटला}, \emph{to āpaṭlā}.

What gets hit takes -\mr{ला} or -\mr{वर}: \mr{दाराला} \mr{आपटणे}, to bang against the door.

\begin{cousinweb}
\textbf{\mr{आपटणे}} (\emph{āpaṭṇe}) is inherited: CDIAL lists Marathi \emph{apaṭṇẽ}, \textquotedblleft{}to strike against\textquotedblright{}, under Sanskrit \emph{āpāṭyate} (CDIAL 1203).
\end{cousinweb}

\begin{grammarlens}[title={What you've built}]
One more word for getting around.
\end{grammarlens}

\subsection*{Your turn}
\begin{itemize}
\raggedright
  \item \emph{Say it:} \emph{āpaṭṇe}
  \item \emph{Say it:} \emph{āpaṭṇe}, once more
  \item \emph{Say it:} \emph{khaḍḍyāmuḷe mājhaṁ ḍokaṁ khiḍkīlā āpaṭlaṁ.} — Because of the pothole my head banged against the window.
  \item \emph{Recall:} say the Marathi for bumpy, then the Marathi for to tip over, then say \emph{āpaṭṇe} again
\end{itemize}

\begin{tcolorbox}[breakable,colback=teal!4,colframe=teal!35!black,title={Before you move on}]
What does \mr{आपटणे} mean? (\textquotedblleft{}To bang against\textquotedblright{}.) Say it once more.
\end{tcolorbox}
\section[(dialogue)]{Looking Back on the Road}

\label{lesson:MR-R403-review-396-403}



\begin{quote}
Say the Marathi for a train and a rail.
\end{quote}

\subsection*{Your turn}
Say these aloud:

\begin{itemize}
\raggedright
  \item a train, a rail, a railway line, a goods train, to rattle
  \item a delay, to be held up, to be kept waiting, to be stalled, to set off
  \item a ticket office, a travel pass, a refund, a destination, to set aside
  \item a driver, a bus conductor, to drive, a fellow traveller, a wayfarer
  \item a vehicle, a two-wheeler, an autorickshaw, a truck, to stuff in
  \item fuel, petrol, a tyre, a puncture, to run out
  \item a ghat road, a narrow pass, a downhill slope, winding, a rock pinnacle
  \item a pothole, a speed breaker, bumpy, to tip over, to bang against
\end{itemize}

\begin{tcolorbox}[breakable,colback=teal!4,colframe=teal!35!black,title={Before you move on}]
A train? (\textbf{\mr{आगगाडी}}, \emph{āggāḍī}.) A rail? (\textbf{\mr{रूळ}}, \emph{rūḷ}.) A railway line? (\textbf{\mr{लोहमार्ग}}, \emph{lohmārga}.) A goods train? (\textbf{\mr{मालगाडी}}, \emph{mālgāḍī}.) To rattle? (\textbf{\mr{खडखडणे}}, \emph{khaḍkhaḍṇe}.) A delay? (\textbf{\mr{विलंब}}, \emph{vilamb}.) To be held up? (\textbf{\mr{खोळंबणे}}, \emph{khoḷambṇe}.) To be kept waiting? (\textbf{\mr{ताटकळणे}}, \emph{tāṭkaḷṇe}.) To be stalled? (\textbf{\mr{रखडणे}}, \emph{rakhaḍṇe}.) To set off? (\textbf{\mr{रवाना} \mr{होणे}}, \emph{ravānā hoṇe}.) A ticket office? (\textbf{\mr{तिकीटघर}}, \emph{tikīṭghar}.) A travel pass? (\textbf{\mr{पास}}, \emph{pās}.) A refund? (\textbf{\mr{परतावा}}, \emph{partāvā}.) A destination? (\textbf{\mr{गंतव्य}}, \emph{gantavya}.) To set aside? (\textbf{\mr{राखून} \mr{ठेवणे}}, \emph{rākhūn ṭhevṇe}.) A driver? (\textbf{\mr{चालक}}, \emph{cālak}.) A bus conductor? (\textbf{\mr{वाहक}}, \emph{vāhak}.) To drive? (\textbf{\mr{चालवणे}}, \emph{cālavṇe}.) A fellow traveller? (\textbf{\mr{सहप्रवासी}}, \emph{sahapravāsī}.) A wayfarer? (\textbf{\mr{वाटसरू}}, \emph{vāṭsarū}.) A vehicle? (\textbf{\mr{वाहन}}, \emph{vāhan}.) A two-wheeler? (\textbf{\mr{दुचाकी}}, \emph{ducākī}.) An autorickshaw? (\textbf{\mr{रिक्षा}}, \emph{rikṣā}.) A truck? (\textbf{\mr{ट्रक}}, \emph{ṭrak}.) To stuff in? (\textbf{\mr{कोंबणे}}, \emph{kombṇe}.) Fuel? (\textbf{\mr{इंधन}}, \emph{indhan}.) Petrol? (\textbf{\mr{पेट्रोल}}, \emph{peṭrol}.) A tyre? (\textbf{\mr{टायर}}, \emph{ṭāyar}.) A puncture? (\textbf{\mr{पंक्चर}}, \emph{paṅkcar}.) To run out? (\textbf{\mr{संपणे}}, \emph{sampṇe}.) A ghat road? (\textbf{\mr{घाट}}, \emph{ghāṭ}.) A narrow pass? (\textbf{\mr{खिंड}}, \emph{khiṇḍ}.) A downhill slope? (\textbf{\mr{उतार}}, \emph{utār}.) Winding? (\textbf{\mr{नागमोडी}}, \emph{nāgmoḍī}.) A rock pinnacle? (\textbf{\mr{सुळका}}, \emph{suḷkā}.) A pothole? (\textbf{\mr{खड्डा}}, \emph{khaḍḍā}.) A speed breaker? (\textbf{\mr{गतिरोधक}}, \emph{gatirodhak}.) Bumpy? (\textbf{\mr{खडबडीत}}, \emph{khaḍbaḍīt}.) To tip over? (\textbf{\mr{कलंडणे}}, \emph{kalaṇḍṇe}.) To bang against? (\textbf{\mr{आपटणे}}, \emph{āpaṭṇe}.)
\end{tcolorbox}
